\chapter{计数器——统一框架变式详解}

\section{变式1：带使能的模N计数器}

\begin{detailbox}[完整教学]
\textbf{【电路】}在标准计数器上增加en控制。
\textbf{【VHDL】}
\begin{lstlisting}
if rising_edge(clk) then
  if en = '1' then
    if cnt = N-1 then cnt <= 0; else cnt <= cnt + 1; end if;
  end if;
end if;
\end{lstlisting}
\textbf{【考点】}en=0时不赋值→保持。这是所有可控计数器的基本范式。
\end{detailbox}


\section{变式2：带同步加载的模N计数器}

\begin{detailbox}[完整教学]
\textbf{【功能】}load=1时cnt=data\_in。load=0时正常计数。
\textbf{【优先级】}load > 计数（if-else中load在前=优先级高）。
\textbf{【VHDL】}
\begin{lstlisting}
if load = '1' then cnt <= data_in;
elsif cnt = N-1 then cnt <= 0;
else cnt <= cnt + 1; end if;
\end{lstlisting}
\end{detailbox}


\section{变式3：递增/递减双向计数器}

\begin{detailbox}[完整教学]
\textbf{【VHDL】}
\begin{lstlisting}
if up = '1' then cnt <= cnt + 1; else cnt <= cnt - 1; end if;
\end{lstlisting}
\textbf{【递减自然溢出】}$0000-1=1111$（4位无符号），自动模16循环。
\end{detailbox}


\section{变式4：倒计时计数器（预置→0→预置）}

\begin{detailbox}[完整教学]
\textbf{【VHDL】}
\begin{lstlisting}
if cnt = 0 then cnt <= PRE;  -- 到0后重新加载
else cnt <= cnt - 1; end if;
\end{lstlisting}
\end{detailbox}


\section{变式5：级联计数器——总模值=乘积}

\begin{detailbox}[完整教学]
\textbf{【级联方式】}低位carry→高位en。总M=$M_1 \times M_2$。

\textbf{【实例】}模8+模5级联→总模40。模60(6×10)BCD级联。

\textbf{【VHDL关键】}高位以低位carry为时钟使能——确保高位只在低位满一轮时才+1。
\end{detailbox}


\section{变式6：格雷码计数器}

\begin{detailbox}[完整教学]
\textbf{【特点】}相邻状态仅1位不同。适合FIFO指针、跨时钟域。
\textbf{【VHDL】}case枚举所有格雷码状态。
\textbf{【3位格雷码序列】}000-001-011-010-110-111-101-100-000。
\end{detailbox}
